\section*{Capítulo \ref{ch:g8:hormonal-communication} --- A comunicação hormonal}

\begin{solution}{exo:g8:hormonal-communication:1}
Liberar o produto dela no \omterm{prop:g4:journey-of-food:absorb}{sangue}, e não por um canal. A \omterm{def:g8:hormonal-communication:glands}{glândula} de
comando da base do \omterm{def:g5:brain-in-command:system}{cérebro}: regente das outras e quem dá a partida na
\omterm{def:g8:puberty:puberty}{puberdade}. A tireoide: a marcadora de ritmo do corpo. As suprarrenais:
o \omterm{def:g8:hormonal-communication:glands}{hormônio} do alarme. O pâncreas: a administração da \omterm{prop:g7:organs-at-work:bill}{glicose} do
\omterm{prop:g4:journey-of-food:absorb}{sangue}. (\omterm{def:g8:reproductive-systems:female}{Ovários} e \omterm{def:g8:reproductive-systems:male}{testículos}: os \omterm{def:g8:hormonal-communication:glands}{hormônios} do ciclo e da
reconstrução.)
\end{solution}

\begin{solution}{exo:g8:hormonal-communication:2}
O transporte é por transmissão aberta, o endereçamento é do lado de
quem recebe: só as \omterm{def:g5:first-look-at-cells:cell}{células} equipadas com os receptores do \omterm{def:g8:hormonal-communication:glands}{hormônio}
respondem. Toda casa recebe a carta; só a chave que encaixa a abre.
\end{solution}

\begin{solution}{exo:g8:hormonal-communication:3}
Ritmo do \omterm{def:g4:heart-and-blood:heart}{coração} para cima (\cref{ch:g7:heart-and-circulation}); vias
aéreas alargadas (\cref{ch:g7:lungs-and-gas-exchange}); o banco de
\omterm{prop:g7:organs-at-work:bill}{glicose} do \omterm{prop:g7:digestion-and-nutrients:absorption}{fígado} liberado
(\cref{ch:g7:digestion-and-nutrients}); \omterm{prop:g4:journey-of-food:absorb}{sangue} redirecionado para os
\omterm{def:g3:bones-and-muscles:muscle}{músculos} (\cref{ch:g7:organs-at-work}) --- com a \omterm{def:g1:the-five-senses:sense}{pele} pálida e
formigando completando o conjunto.
\end{solution}

\begin{solution}{exo:g8:hormonal-communication:4}
A \omterm{def:g8:hormonal-communication:glands}{glândula} que escreve monitora o resultado das cartas dela e diminui
o ritmo conforme a meta é alcançada --- termostato, e não torneira.
Vitrine: a \omterm{prop:g7:organs-at-work:bill}{glicose} do \omterm{prop:g4:journey-of-food:absorb}{sangue}, depositada na subida e liberada na
queda pelos \omterm{def:g8:hormonal-communication:glands}{hormônios} em par do pâncreas.
\end{solution}

\begin{solution}{exo:g8:hormonal-communication:5}
O sinal do \omterm{def:g8:hormonal-communication:glands}{hormônio} de estoque --- não escrito ou não respondido pelos
receptores dele. A \omterm{prop:g7:organs-at-work:bill}{glicose} então se acumula depois das refeições,
sobrecarrega a \omterm{def:g7:organs-at-work:recovery}{recuperação} do \omterm{def:g7:removing-wastes:kidneys}{rim} e aparece na \omterm{def:g7:removing-wastes:kidneys}{urina}.
\end{solution}

\begin{solution}{exo:g8:hormonal-communication:6}
Na base do \omterm{def:g5:brain-in-command:system}{cérebro} --- a \omterm{def:g8:hormonal-communication:glands}{glândula} de comando, com o \omterm{def:g5:brain-in-command:system}{sistema nervoso}
acima e os \omterm{def:g8:hormonal-communication:glands}{hormônios} abaixo. Exemplo: o susto --- primeiro os sinais
nervosos, com a carta das suprarrenais sustentando o alerta; ou a
duração do dia virando os \omterm{def:g8:hormonal-communication:glands}{hormônios} de reprodução das ovelhas.
\end{solution}

\begin{solution}{exo:g8:hormonal-communication:7}
Uma torneira despeja até alguém fechar; um termostato mede aquilo que
ele está mudando e se ajusta sozinho. A \omterm{def:g8:hormonal-communication:glands}{glândula} lê o nível que ela
regula --- a secreção sobe e diminui junto com a necessidade, sem
precisar de nenhuma mão de fora.
\end{solution}

\begin{solution}{exo:g8:hormonal-communication:8}
Quem escreve: o pâncreas, disparado pela \omterm{prop:g7:organs-at-work:bill}{glicose} do \omterm{prop:g4:journey-of-food:absorb}{sangue} que sobe.
Quem recebe: o \omterm{prop:g7:digestion-and-nutrients:absorption}{fígado} e os \omterm{def:g3:bones-and-muscles:muscle}{músculos} acima de tudo --- eles depositam a
sobra ao receber. \omterm{prop:g8:hormonal-communication:feedback}{Retroalimentação}: o próprio nível que cai faz a
secreção diminuir. Falha e uso: o \omterm{ex:g8:hormonal-communication:diabetes}{diabetes} --- e as doses medidas da
medicina, a carta por seringa.
\end{solution}

\begin{solution}{exo:g8:hormonal-communication:9}
Os dois escrevem à mão uma carta do corpo: a seringa fornece a
mensagem de estoque que falta ao pâncreas na hora das refeições; a
pílula fornece a mensagem de ``suspenso'' do ciclo todo dia. O mesmo
ato --- imitar um \omterm{def:g8:hormonal-communication:glands}{hormônio} para rodar o sistema dele de propósito ---
sendo que um restaura uma correspondência que falhou e o outro
sobrepõe-se por escolha a uma correspondência saudável.
\end{solution}

\begin{solution}{exo:g8:hormonal-communication:10}
O início do tremor de frio: nervoso --- ordens musculares em fração de
segundo. A barba: hormonal --- meses, corpo inteiro, corrente
sanguínea. A faixa da \omterm{prop:g7:organs-at-work:bill}{glicose}: hormonal --- correspondência com
\omterm{prop:g8:hormonal-communication:feedback}{retroalimentação}. O reflexo patelar: nervoso --- uma junção, um
cinquentavo de segundo.
\end{solution}

\begin{solution}{exo:g8:hormonal-communication:11}
A percepção dos dias que encurtam, pelos \omterm{def:g1:the-five-senses:sense}{olhos}, é nervosa; o
escritório da base do \omterm{def:g5:brain-in-command:system}{cérebro} traduz essa percepção e rege a \omterm{def:g8:hormonal-communication:glands}{glândula}
de comando; as cartas dela acordam o programa hormonal dos \omterm{def:g8:reproductive-systems:female}{ovários}; os
ciclos das ovelhas começam --- \omterm{def:g6:exploring-our-environment:conditions}{luz} entrando pelo telégrafo, cordeiros
saindo pelo correio.
\end{solution}

\begin{solution}{exo:g8:hormonal-communication:12}
Defesa: um \omterm{prop:g4:journey-of-food:absorb}{sangue} só carrega todas as cartas e, mesmo assim, cada uma
age só onde há receptores que respondem --- a plateia é que define o
canal, e assim a transmissão aberta vira linhas privadas. O passo
adiante: uma \omterm{def:g5:first-look-at-cells:cell}{célula} entra numa plateia construindo o receptor --- então
que tecidos escutam é algo construído, órgão por órgão, durante o
desenvolvimento; os alvos da \omterm{def:g8:puberty:puberty}{puberdade} --- laringe, \omterm{def:g1:the-five-senses:sense}{pele}, placas de
crescimento --- foram equipados com os receptores deles de antemão,
esperando anos pela primeira carta.
\end{solution}

\begin{solution}{pb:g8:hormonal-communication:1}
\textbf{1.} As cartas em par do pâncreas --- o \omterm{def:g8:hormonal-communication:glands}{hormônio} de estoque na
subida e o parceiro que libera na queda --- segurando a faixa por
\omterm{prop:g8:hormonal-communication:feedback}{retroalimentação} enquanto o dono dormia.

\textbf{2.} Disparador: a \omterm{prop:g7:organs-at-work:bill}{glicose} da refeição atravessando as
\omterm{prop:g7:digestion-and-nutrients:absorption}{vilosidades} e levantando o nível do \omterm{prop:g4:journey-of-food:absorb}{sangue}. Quem escreve: o pâncreas.
Carta: o \omterm{def:g8:hormonal-communication:glands}{hormônio} de estoque. Receptores: o \omterm{prop:g7:digestion-and-nutrients:absorption}{fígado} e os \omterm{def:g3:bones-and-muscles:muscle}{músculos},
depositando a sobra.

\textbf{3.} Uma torneira despejaria além da meta, até o mergulho; a
descida se aplainando dentro da faixa mostra a secreção diminuindo à
medida que o próprio resultado dela chega --- automedição e
autoparada.

\textbf{4.} Que ela subiu junto com o nível e recuou junto com ele:
uma curva de secreção acompanhando a curva da \omterm{prop:g7:organs-at-work:bill}{glicose}, com o pico
logo atrás do pico --- o percurso do termostato.

\textbf{5.} O parceiro que libera (com a ajuda do \omterm{def:g8:hormonal-communication:glands}{hormônio} do alarme
no esforço pesado): a \omterm{prop:g7:organs-at-work:bill}{glicose} depositada no \omterm{prop:g7:digestion-and-nutrients:absorption}{fígado} --- o banco enchido
no café da manhã, à maneira do
\cref{ex:g7:digestion-and-nutrients:breadtrace} --- gasta na conta dos
\omterm{def:g3:bones-and-muscles:muscle}{músculos} em atividade.

\textbf{6.} O mesmo parceiro que libera: destrancando o banco de forma
constante ao longo das horas sem comida, acompanhando o gasto do corpo
--- a faixa segurada pelo outro lado.

\textbf{7.} As \omterm{def:g8:hormonal-communication:glands}{glândulas} suprarrenais. O açúcar prepara combustível
para \emph{fugir ou enfrentar}: o objetivo inteiro da carta do alarme
é um corpo abastecido e redirecionado para um esforço imediato ---
venha ou não o esforço.

\textbf{8.} Porque a correspondência saudável respondeu: a sobra
liberada, não gasta (o quase acidente passou), disparou por sua vez a
carta de estoque, e a faixa foi restaurada --- alarme e termostato em
sequência.

\textbf{9.} Sem tratamento: a subida do café da manhã continuando a
subir, sem descer --- nenhuma carta de estoque escrita ou escutada ---
um platô alto durante a manhã inteira; passado o limiar da
\omterm{def:g7:organs-at-work:recovery}{recuperação}, aparece açúcar na \omterm{def:g7:removing-wastes:kidneys}{urina} (o sinal do
\cref{ex:g7:removing-wastes:thrift}).

\textbf{10.} A carta que o próprio pâncreas escreveria na hora da
refeição. A dose precisa combinar com o prato porque o tamanho da
carta precisa combinar com a \omterm{prop:g7:organs-at-work:bill}{glicose} que ela vai depositar --- pequena
demais deixa um acúmulo, grande demais ultrapassa a faixa para baixo,
o que é uma emergência por si só.

\textbf{11.} Noite e tarde: o par do pâncreas, com o \omterm{prop:g7:digestion-and-nutrients:absorption}{fígado} e os
\omterm{def:g3:bones-and-muscles:muscle}{músculos} respondendo --- a faixa segurada. Refeições da manhã e da
noite: a \omterm{prop:g7:organs-at-work:bill}{glicose} fala, o pâncreas escreve, o \omterm{prop:g7:digestion-and-nutrients:absorption}{fígado} deposita.
Meio-dia: o esforço gasta, o parceiro que libera destranca. Cinco da
tarde: as suprarrenais interrompem, e o pâncreas arruma depois. Um
valor, cinco falantes, nenhum silêncio o dia inteiro.

\textbf{12.} ``Um sistema de \omterm{prop:g8:hormonal-communication:feedback}{retroalimentação}: escritoras que leem os
resultados das próprias cartas --- é assim que um valor segue por uma
faixa estreita a vida inteira, sem ninguém dirigindo.''
\end{solution}
